\documentclass[11pt,a4paper]{article}

\usepackage[margin=2.2cm]{geometry}
\usepackage[T1]{fontenc}
\usepackage{lmodern}
\usepackage{booktabs}
\usepackage{tabularx}
\usepackage{xcolor}
\usepackage{listings}
\usepackage{enumitem}
\usepackage[hidelinks]{hyperref}

\lstset{
  basicstyle=\ttfamily\footnotesize,
  columns=fullflexible,
  keepspaces=true,
  frame=single,
  rulecolor=\color{gray!60},
  backgroundcolor=\color{gray!6},
  breaklines=true,
  showstringspaces=false,
  xleftmargin=2pt,
  xrightmargin=2pt
}
\lstdefinestyle{py}{
  language=Python,
  keywordstyle=\color{blue!60!black}\bfseries,
  commentstyle=\color{green!40!black}\itshape,
  stringstyle=\color{red!50!black}
}

\setlist{itemsep=2pt, topsep=4pt}
\setlength{\parskip}{4pt}
\setlength{\parindent}{0pt}

\title{\textbf{Programming Assignment 3: Mini Packet Analyzer}\\[2pt]
\large Computer Networks}
\author{Name: \underline{\hspace{5cm}} \quad Roll No.: \underline{\hspace{3cm}}}
\date{}

\begin{document}
\maketitle
\vspace{-1.5em}

\section{Overview}

The goal of this assignment was to build a small, command-line version of
Wireshark in Python. The tool captures IPv4 packets sent or received by the
host, numbers and timestamps them, and decodes the IPv4, TCP and UDP headers.
Captured packets can be listed, filtered and opened one at a time to see the
full header breakdown and the payload.

Only the Python standard library is used (\texttt{socket}, \texttt{struct},
\texttt{ipaddress}, \texttt{datetime}, \texttt{argparse}). Sockets are used
just to get the raw bytes; every header field is pulled out of those bytes by
our own code. No Scapy, PyShark or similar decoding library is involved.

\section{Program Structure}

The code is split into small modules so that capturing, parsing and printing
stay independent of each other.

\begin{center}
\begin{tabularx}{\textwidth}{@{}lX@{}}
\toprule
\textbf{File} & \textbf{Responsibility} \\
\midrule
\texttt{packet\_capture.py} & Opens a raw socket (Windows/Linux) and yields raw IPv4 bytes with a timestamp. \\
\texttt{packet\_parser.py} & Manual parsing of IPv4, TCP and UDP headers; the core of the assignment. \\
\texttt{packet\_display.py} & Summary table, filters, detailed view and hex/ASCII payload dump. \\
\texttt{mini\_packet\_analyzer.py} & Command-line options and the interactive \texttt{analyzer>} prompt. \\
\texttt{demo\_packets.py} & Builds seven hand-made packets so the tool can be shown without admin rights. \\
\texttt{test\_analyzer.py} & 31 unit tests. \\
\bottomrule
\end{tabularx}
\end{center}

The flow is simple: \emph{capture} $\rightarrow$ \emph{parse} $\rightarrow$
\emph{store as a numbered Packet} $\rightarrow$ \emph{filter summary}
$\rightarrow$ \emph{select a packet and show details}.

\section{Packet Capture}

\begin{itemize}
  \item \textbf{Linux:} an \texttt{AF\_PACKET} raw socket with
  \texttt{ETH\_P\_ALL} receives both incoming and outgoing frames. These are
  Ethernet frames, so the 14-byte Ethernet header (and any 802.1Q VLAN tags)
  is stripped first, and only frames with EtherType \texttt{0x0800} (IPv4) are
  kept. On the loopback interface Linux delivers every packet twice, so the
  outgoing copy is dropped.
  \item \textbf{Windows:} an \texttt{AF\_INET} raw socket is bound to the
  adapter's IPv4 address and switched to receive-all mode with
  \texttt{SIO\_RCVALL}. Windows already hands over IP packets, so nothing has
  to be stripped. Receive-all mode is turned off again when capture ends.
\end{itemize}

Capture stops after \texttt{--count} packets, after \texttt{--timeout}
seconds, or on Ctrl+C (the packets collected so far are kept). Each packet gets
a sequential number and a timestamp taken the moment it is received. The size
shown is the number of IPv4 bytes (Ethernet header excluded). Root/CAP\_NET\_RAW
on Linux or Administrator on Windows is needed; without it the program prints a
clear error. A \texttt{--demo} mode feeds synthetic packets through the exact
same parser for testing without privileges.

\section{Header Parsing}

All multi-byte fields are big-endian (network byte order).
\texttt{struct.unpack} with the \texttt{!} prefix is used to turn fixed-size
slices into integers; it is only a byte converter and knows nothing about
protocols, so the offsets, sizes and bit masks are all chosen by us.

\subsection{IPv4}

\begin{center}
\begin{tabular}{@{}cll@{}}
\toprule
\textbf{Bytes} & \textbf{Field} & \textbf{How it is extracted} \\
\midrule
0      & Version, IHL        & \texttt{raw[0] >> 4}, \;\texttt{(raw[0] \& 0x0F) * 4} \\
2--3   & Total length        & 16-bit unsigned (\texttt{H}) \\
4--5   & Identification      & \texttt{H} \\
6--7   & Flags, frag. offset & DF \texttt{\& 0x4000}, MF \texttt{\& 0x2000}, offset \texttt{\& 0x1FFF} \\
8      & TTL                 & 1 byte (\texttt{B}) \\
9      & Protocol            & \texttt{B}; 6 = TCP, 17 = UDP \\
10--11 & Header checksum     & \texttt{H} \\
12--15 & Source IP           & 4 bytes, \texttt{socket.inet\_ntoa} \\
16--19 & Destination IP      & 4 bytes, \texttt{socket.inet\_ntoa} \\
\bottomrule
\end{tabular}
\end{center}

The IHL is a count of 32-bit words, which is why it is multiplied by 4. This
matters because the IP header can carry options; the transport header starts at
\texttt{raw[header\_length:]}, not blindly at byte 20.

\begin{lstlisting}[style=py]
version       = raw[0] >> 4
header_length = (raw[0] & 0x0F) * 4
fields = struct.unpack("!BBHHHBBH4s4s", raw[:20])
ttl, protocol, checksum = fields[5], fields[6], fields[7]
source = socket.inet_ntoa(fields[8])
\end{lstlisting}

The protocol byte is what lets the program decide automatically whether the
packet carries TCP or UDP. Anything else (e.g.\ ICMP) is shown with its IP
header only.

\subsection{TCP}

\begin{center}
\begin{tabular}{@{}cll@{}}
\toprule
\textbf{Bytes} & \textbf{Field} & \textbf{How it is extracted} \\
\midrule
0--1   & Source port         & \texttt{H} \\
2--3   & Destination port    & \texttt{H} \\
4--7   & Sequence number     & 32-bit unsigned (\texttt{I}) \\
8--11  & Acknowledgment no.  & \texttt{I} \\
12     & Data offset         & \texttt{(byte >> 4) * 4} \\
13     & Flags               & bit 0 FIN, 1 SYN, 2 RST, 3 PSH, 4 ACK, 5 URG, 6 ECE, 7 CWR \\
14--15 & Window size         & \texttt{H} \\
16--17 & Checksum            & \texttt{H} \\
18--19 & Urgent pointer      & \texttt{H} \\
\bottomrule
\end{tabular}
\end{center}

\begin{lstlisting}[style=py]
fields = struct.unpack("!HHIIBBHHH", segment[:20])
header_length = (fields[4] >> 4) * 4
names = ["FIN","SYN","RST","PSH","ACK","URG","ECE","CWR"]
flags = [n for bit, n in enumerate(names) if fields[5] & (1 << bit)]
payload_length = len(segment) - header_length
\end{lstlisting}

\subsection{UDP}

UDP has a fixed 8-byte header: source port, destination port, length and
checksum, each 2 bytes (\texttt{"!HHHH"}). The length field covers header plus
data, so the payload is \texttt{segment[8:length]}, and its size is
\texttt{length - 8}.

\subsection{Payload and edge cases}

The payload is whatever follows the transport header, up to the IP total
length. A few cases needed extra care:

\begin{itemize}
  \item \textbf{Empty payload:} a pure ACK or SYN shows \texttt{Payload length:
  0} and \texttt{(empty)} rather than trying to decode non-existent data.
  \item \textbf{Ethernet padding:} small frames are padded to 60 bytes; anything
  past the IP total length is cut off so it isn't mistaken for data.
  \item \textbf{Total length = 0:} outgoing packets can show this because of
  segmentation offload on the NIC; the captured length is used instead.
  \item \textbf{Fragments and malformed packets:} later IP fragments are not
  parsed as TCP/UDP (they start mid-segment). Truncated or invalid headers are
  kept in the list with a note instead of crashing the program.
\end{itemize}

The payload is printed as a hex dump with a printable-ASCII column, like
Wireshark. Non-printable bytes appear as \texttt{.}, so binary or encrypted
data can't mess up the terminal.

\section{Summary, Filtering and Selection}

After capture, a summary table lists the packet number, time, source and
destination IP, protocol, source and destination port, and length. Filters can
be given on the command line or typed at the interactive prompt:

\begin{lstlisting}
--protocol TCP|UDP   --src-ip A.B.C.D   --dst-ip A.B.C.D
--src-port N         --dst-port N
\end{lstlisting}

When several filters are given they are combined with AND. Filtering only hides
rows; packets keep their original numbers, and \texttt{reset} shows everything
again. Typing a packet number at the \texttt{analyzer>} prompt (or passing
\texttt{--packet N}) prints the detailed view.

\section{Sample Output}

Summary with a UDP + destination-port filter applied:

\begin{lstlisting}
$ python mini_packet_analyzer.py --demo --protocol UDP --dst-port 9999
  No.  Time          Source IP     Destination IP  Protocol  Src port  Dst port  IP bytes
    3  12:27:33.130  192.168.1.10  192.168.1.20    UDP          53000      9999        43
    7  12:27:33.530  192.168.1.10  192.168.1.20    UDP          53001      9999        35
Shown: 2 / Stored: 7.
\end{lstlisting}

Detailed view of a TCP packet:

\begin{lstlisting}
Packet #1
Time:                    2026-10-29 12:27:32.930+05:30
Packet length:           52 IPv4 bytes (without Ethernet)

IP Header
------------------------------------------------
Version:                 4
Header length:           20 bytes
Total length:            52 bytes
TTL:                     64
Protocol:                TCP (6)
Header checksum:         0xB283
Source IP:               192.168.1.10
Destination IP:          192.168.1.20

TCP Header
------------------------------------------------
Source port:             49152
Destination port:        8080
Sequence number:         105482
Acknowledgment number:   208713
Header length:           20 bytes
Flags:                   PSH, ACK (0x18)
Window size:             65535
Checksum:                0xB2C9
Payload length:          12 bytes available

Payload: hex bytes | printable ASCII
------------------------------------------------
0000  48 45 4C 4C 4F 20 53 45 52 56 45 52              |HELLO SERVER|
\end{lstlisting}

A SYN packet with IP and TCP options (demo packet 6) correctly reports 24-byte
IP and TCP headers, shows the option bytes in hex, and prints an empty payload.

\section{Testing}

\begin{itemize}
  \item \textbf{Unit tests:} 31 tests (\texttt{python -m unittest
  test\_analyzer.py}), all passing. They check every header field against
  hand-built byte strings, all eight TCP flags, IP/TCP options, empty payloads,
  each filter alone and combined, truncated/malformed input, fragments,
  Ethernet/VLAN stripping and socket clean-up (using mocked sockets).
  \item \textbf{Live capture (Linux, root):} a UDP datagram with the text
  \texttt{PINGTEST} sent to \texttt{127.0.0.1:9999} was captured on
  \texttt{lo} and decoded correctly. An HTTPS request captured on
  \texttt{eth0} showed both outgoing and incoming TCP packets on port 443.
  \item The demo packets have valid IP checksums, verified in the tests.
\end{itemize}

\section{Limitations}

\begin{itemize}
  \item Only IPv4 is decoded; IPv6 and ARP frames are skipped.
  \item Checksums are displayed but not verified (checksum offload makes
  outgoing values unreliable anyway).
  \item IP fragments are not reassembled.
  \item Packets are kept in memory only; there is no pcap export.
  \item On Windows, one adapter is captured per run, and loopback traffic may
  not be visible.
\end{itemize}

\section{Conclusion}

Writing the parser by hand made it clear how a packet is just a sequence of
bytes whose meaning comes entirely from fixed offsets defined by the protocol.
Small details, such as the header length being in 4-byte words, flags packed
into single bits, and the need to respect the IP total length, turned out to be
exactly where a naive parser goes wrong. The finished tool covers capture,
per-field decoding of IPv4/TCP/UDP, filtering and per-packet inspection as
required.

\end{document}
